\documentclass[UTF8,fontset=none,12pt]{ctexart}

\setCJKmainfont{Songti SC}
\setCJKsansfont{PingFang SC}
\setCJKmonofont{Songti SC}

\usepackage[a4paper,margin=2.35cm]{geometry}
\usepackage{amsmath,amssymb,amsthm,mathtools,bm}
\usepackage{graphicx}
\usepackage{float}
\usepackage{booktabs}
\usepackage{array}
\usepackage{enumitem}
\usepackage[hidelinks]{hyperref}

\title{数值线性代数\quad 第 29 讲：移位 QR 算法}
\author{Lloyd N. Trefethen、David Bau III\quad 中文阅读版}
\date{}

\begin{document}
\maketitle
\pagestyle{plain}

让 QR 迭代真正“飞起来”的关键，是在每一步引入移位
$A\mathbin{\mapsto}A-\mu I$。本讲解释这一思想如何通过与 Rayleigh 商迭代
的隐含联系带来三次收敛。

\section*{与逆迭代的联系}

继续假设 $A\in\mathbb{R}^{m\times m}$ 是实对称矩阵，特征值为实数
$\{\lambda_j\}$，特征向量 $\{q_j\}$ 标准正交。

我们已经看到，“纯” QR 算法（算法 28.1）等价于把同时迭代应用于单位矩阵；
特别地，结果的第一列按照从 $e_1$ 出发的幂迭代演化。这个观察有一个对偶版本：
算法 28.1 也等价于把同时逆迭代应用于一个“翻转后的”单位矩阵 $P$，因此结果
的第 $m$ 列按照从 $e_m$ 出发的逆迭代演化。

下面建立这个结论。为避免原书的下划线记号在纯文本中产生歧义，本文用
$Q^{(j)}$ 表示第 $j$ 步 QR 分解得到的正交因子，并用

$$
\mathcal Q^{(k)}=Q^{(1)}Q^{(2)}\cdots Q^{(k)}
$$

表示累计的正交因子；相应的累计上三角因子记为 $\mathcal R^{(k)}$。上一讲的
式 (28.14) 说明

$$
\mathcal Q^{(k)}=\prod_{j=1}^{k}Q^{(j)}
 =\begin{bmatrix}q_1^{(k)}&q_2^{(k)}&\cdots&q_m^{(k)}\end{bmatrix},
$$

它正是同时迭代在第 $k$ 步（式 (28.9)）出现的正交矩阵。换一种说法，
$\mathcal Q^{(k)}$ 是 QR 分解

$$
A^k=\mathcal Q^{(k)}\mathcal R^{(k)}
$$

中的正交因子。

现在把这个公式求逆。由于 $A^{-1}$ 仍然是对称矩阵，有

\begin{equation}
A^{-k}=\bigl(\mathcal R^{(k)}\bigr)^{-1}
          \bigl(\mathcal Q^{(k)}\bigr)^T
        =\mathcal Q^{(k)}\bigl(\mathcal R^{(k)}\bigr)^{-T}.
\tag{29.1}
\end{equation}

令 $P$ 为把行或列的顺序完全翻转的 $m\times m$ 置换矩阵：

$$
P=\begin{bmatrix}
0&\cdots&0&1\\
\vdots&\ddots&1&0\\
0&1&\ddots&\vdots\\
1&0&\cdots&0
\end{bmatrix},
\qquad P^2=I.
$$

于是 (29.1) 可以改写为

\begin{equation}
A^{-k}P
 =\bigl[\mathcal Q^{(k)}P\bigr]
  \bigl[P\bigl(\mathcal R^{(k)}\bigr)^{-T}P\bigr].
\tag{29.2}
\end{equation}

乘积中的第一个因子 $\mathcal Q^{(k)}P$ 是正交矩阵。第二个因子
$P(\mathcal R^{(k)})^{-T}P$ 是上三角矩阵：先从下三角矩阵
$(\mathcal R^{(k)})^{-T}$ 出发，上下翻转一次，再左右翻转一次即可看到这一点。
因此，(29.2) 可以解释为 $A^{-k}P$ 的 QR 分解。换句话说，我们实际上是在
对 $A^{-1}$ 做初始矩阵为 $P$ 的同时迭代，也就是对 $A$ 做同时逆迭代。
特别地，$\mathcal Q^{(k)}P$ 的第一列——也就是 $\mathcal Q^{(k)}$ 的最后一列——
正是对向量 $e_m$ 做 $k$ 步逆迭代所得的结果。

\section*{与移位逆迭代的联系}

因此，QR 算法同时具有同时迭代和同时逆迭代两种解释；这两个视角完全对称。
但是，正如第 27 讲所见，幂迭代和逆迭代之间有一个巨大差异：逆迭代可以通过
移位任意加速。我们越能准确估计某个特征值 $\mu\approx\lambda_j$，用移位矩阵
$A-\mu I$ 做一步逆迭代时得到的效果就越好。算法 28.2 说明了如何在 QR 的一步
迭代中引入移位；这恰好对应于同时迭代和同时逆迭代中的移位，因此它们的有利
效果也完全相同。

令 $\mu^{(k)}$ 表示 QR 算法第 $k$ 步选取的特征值估计。由算法 28.2，从第
$k-1$ 步到第 $k$ 步的移位 QR 关系为

$$
A^{(k-1)}-\mu^{(k)}I=Q^{(k)}R^{(k)},
\qquad
A^{(k)}=R^{(k)}Q^{(k)}+\mu^{(k)}I.
$$

这意味着

\begin{equation}
A^{(k)}=(Q^{(k)})^TA^{(k-1)}Q^{(k)},
\tag{29.3}
\end{equation}

并且通过归纳得到

$$
A^{(k)}=(\mathcal Q^{(k)})^TA\mathcal Q^{(k)},
$$

这与式 (28.17) 没有变化。不过，式 (28.16) 不再成立。现在得到的分解是

\begin{equation}
\begin{aligned}
&(A-\mu^{(k)}I)(A-\mu^{(k-1)}I)\cdots(A-\mu^{(1)}I)\\
&\hspace{4cm}=\mathcal Q^{(k)}\mathcal R^{(k)}.
\end{aligned}
\tag{29.4}
\end{equation}

这是同时迭代的一个移位版本（证明略）。换句话说，
$\mathcal Q^{(k)}=\prod_{j=1}^kQ^{(j)}$ 是乘积
$\prod_{j=1}^k(A-\mu^{(j)}I)$ 的正交化结果。$\mathcal Q^{(k)}$ 的第一列是
使用这些移位从 $e_1$ 出发的移位幂迭代结果；最后一列则是使用同一组移位从
$e_m$ 出发做 $k$ 步移位逆迭代的结果。如果移位是足够好的特征值估计，最后一列
就会快速收敛到一个特征向量。

\section*{与 Rayleigh 商迭代的联系}

我们已经在移位 QR 算法中发现了一个强大的工具：移位逆迭代。要使 QR 算法最后
一列快速收敛，还需要解决如何选择移位的问题。

Rayleigh 商是一个自然的起点。若要估计 $\mathcal Q^{(k)}$ 最后一列所近似的
特征向量对应的特征值，最自然的做法就是对这一列使用 Rayleigh 商。这给出

\begin{equation}
\mu^{(k)}
 =\frac{(q_m^{(k)})^TAq_m^{(k)}}{(q_m^{(k)})^Tq_m^{(k)}}
 =(q_m^{(k)})^TAq_m^{(k)}.
\tag{29.5}
\end{equation}

如果每一步都把这个数选作移位，那么特征值估计 $\mu^{(k)}$ 和特征向量估计
$q_m^{(k)}$ 就与从 $e_m$ 出发执行 Rayleigh 商迭代所得到的结果完全相同。
因此，QR 算法具有三次收敛性，具体说是 $q_m^{(k)}$ 三次收敛到一个特征向量。

注意，在 QR 算法中，$q_m^{(k)}$ 的 Rayleigh 商就是 $A^{(k)}$ 的 $(m,m)$
条目，因此这个移位几乎不需要额外计算。上一讲结尾已经提到过这一点，这里
再明确推导一次。从 (29.3) 出发，有

\begin{equation}
A_{mm}^{(k)}=e_m^TA^{(k)}e_m
 =e_m^T(\mathcal Q^{(k)})^TA\mathcal Q^{(k)}e_m
 =(q_m^{(k)})^TAq_m^{(k)}.
\tag{29.6}
\end{equation}

因此，(29.5) 等价于简单地取
$\mu^{(k)}=A_{mm}^{(k)}$。这称为 Rayleigh 商移位。

\section*{Wilkinson 移位}

Rayleigh 商移位在一般情形下给出三次收敛，但并不能保证所有初始条件都收敛。
用一个简单例子即可看出这一点。考虑矩阵

\begin{equation}
A=\begin{bmatrix}0&1\\1&0\end{bmatrix}.
\tag{29.7}
\end{equation}

无移位 QR 算法完全不收敛：

$$
A=Q^{(1)}R^{(1)}
 =\begin{bmatrix}0&1\\1&0\end{bmatrix}
  \begin{bmatrix}1&0\\0&1\end{bmatrix},
$$

而

$$
A^{(1)}=R^{(1)}Q^{(1)}
 =\begin{bmatrix}1&0\\0&1\end{bmatrix}
  \begin{bmatrix}0&1\\1&0\end{bmatrix}=A.
$$

不过，Rayleigh 商移位 $\mu=A_{mm}$ 也没有作用，因为 $A_{mm}=0$。因此在
最坏情形下，带 Rayleigh 商移位的 QR 算法可能失败。

问题来自特征值的对称性：一个特征值为 $+1$，另一个为 $-1$。当我们试图改进
特征值估计 0 时，算法偏向两个特征值的倾向完全相同，因而估计不会改进。我们
需要一种能够打破这种对称性的特征值估计。下面给出一种选择。令 $B$ 表示
$A^{(k)}$ 右下角的 $2\times2$ 子矩阵：

$$
B=\begin{bmatrix}
a_{m-1}&b_{m-1}\\
b_{m-1}&a_m
\end{bmatrix}.
$$

Wilkinson 移位定义为 $B$ 中距离 $a_m$ 更近的那个特征值；如果两个特征值
距离相同，则任意选取其中一个。Wilkinson 移位的一个数值稳定公式是

\begin{equation}
\mu=a_m-\operatorname{sign}(\delta)\,
\frac{b_{m-1}^2}{|\delta|+\sqrt{\delta^2+b_{m-1}^2}},
\qquad
\delta=\frac{a_{m-1}-a_m}{2}.
\tag{29.8}
\end{equation}

如果 $\delta=0$，$\operatorname{sign}(\delta)$ 可以任意设为 $1$ 或 $-1$。

和 Rayleigh 商移位一样，Wilkinson 移位在一般情形下达到三次收敛；此外可以
证明，在最坏情形下它至少达到二次收敛。特别地，带 Wilkinson 移位的 QR 算法
在精确算术中总会收敛。

在例 (29.7) 中，Wilkinson 移位为 $+1$ 或 $-1$。因此对称性被打破，算法一步
就会收敛。

\section*{稳定性与精度}

至此，我们完成了对 QR 算法机制的讨论；许多实际细节仍被省略，例如降阶条件
以及移位的“隐式”实现策略。最后说明一下稳定性和精度。

QR 算法使用正交矩阵，因此它后向稳定并不令人意外。和前面各讲一样，最简单的
表述方式是：令 $\widetilde\Lambda$ 表示浮点运算中计算得到的 $A$ 的对角化，
令 $\widetilde Q$ 表示所有数值计算得到的 Householder 反射（或 Givens 旋转）
之积所对应的、严格正交的矩阵。可以证明如下结果。

\paragraph{定理 29.1.}
设实对称三对角矩阵 $A\in\mathbb{R}^{m\times m}$ 在满足 (13.5)、(13.7)
的计算机上由 QR 算法（算法 28.2）对角化，并按上述方式定义
$\widetilde\Lambda$ 和 $\widetilde Q$。则存在某个
$\delta A\in\mathbb{C}^{m\times m}$，使得

\begin{equation}
\widetilde Q\,\widetilde\Lambda\,\widetilde Q^*
 =A+\delta A,
\qquad
\frac{\|\delta A\|}{\|A\|}=O(\varepsilon_{\mathrm{machine}}).
\tag{29.9}
\end{equation}

和本书中的大多数算法一样，QR 算法给出了一个略有扰动的问题的精确解。结合
定理 26.1 和定理 29.1，可以看到，三对角化之后再执行 QR 算法，是一个计算矩阵
特征值的后向稳定算法。再把这个结论与式 (26.4) 关于实对称矩阵（正规矩阵的
特殊情形）特征值扰动的结果结合起来，就得到计算特征值 $\widetilde\lambda_j$
满足

\begin{equation}
\frac{|\widetilde\lambda_j-\lambda_j|}{\|A\|}
 =O(\varepsilon_{\mathrm{machine}}).
\tag{29.10}
\end{equation}

对于一个只需要约 $\frac{4}{3}m^3$ 次浮点运算的算法来说，这个结果相当不错；
它只相当于计算两个 $m\times m$ 矩阵乘积所需运算量的三分之二。

\section*{习题}

\paragraph{29.1.}
这道五部分习题要求只使用基本构件编写一个 MATLAB 程序，求出实对称矩阵的
全部特征值。不必通过最优地利用对称性或零结构来追求最优常数；整个问题可以
用大约 50 行程序解决。

\begin{enumerate}[label=(\alph*)]
  \item 编写函数 $T=\texttt{tridiag}(A)$，通过正交相似变换把实对称
  $m\times m$ 矩阵约化为三对角形式。程序只能使用 MATLAB 的基本操作，不能
  使用内置函数 `hess` 等。除舍入误差外，输出矩阵 $T$ 应当对称且为三对角矩阵；
  如果愿意，可以在最后加一行强制 $T$ 精确对称且精确三对角。以
  $A=\texttt{hilb}(4)$ 为例运行程序。

  \item 编写函数 $T_{\mathrm{new}}=\texttt{qralg}(T)$，对实三对角矩阵
  $T$ 执行无移位 QR 算法。每一步的 QR 分解可以使用习题 10.2 中的程序
  $[W,R]=\texttt{house}(A)$ 和 $Q=\texttt{formQ}(W)$（如果已经有），也可以
  使用 MATLAB 命令 `qr`；为了提高效率，还可以编写使用 Givens 旋转或
  $2\times2$ Householder 反射的新程序，避免 $m\times m$ 运算。也可以在每一步
  强制对称性和三对角结构。当
  
  $$
  |t_{m,m-1}|<10^{-12}
  $$
  
  时停止，并将当前三对角矩阵 $T$ 作为 $T_{\mathrm{new}}$ 返回（这还远算不上
  工业级的收敛判据）。同样对 $A=\texttt{hilb}(4)$ 运行。

  \item 编写驱动程序：(i) 调用 `tridiag`；(ii) 调用 `qralg` 得到一个
  特征值；(iii) 对更小的矩阵调用 `qralg` 得到另一个特征值；如此继续，直到
  得到 $A$ 的全部特征值。设置程序，使其把每一次 QR 迭代中的
  $|t_{m,m-1}|$ 存入一个向量，并在最后生成这些数值相对于 QR 分解次数的
  `semilogy` 图。随着降阶进行，$m$ 会从 `length(A)` 逐步变为
  `length(A)-1`，一直到 3，最后到 2，因此图形会呈现锯齿状。对
  $A=\texttt{hilb}(4)$ 运行程序；输出应包括一组特征值和一幅“锯齿图”。

  \item 修改 `qralg`，使其在每一步使用 Wilkinson 移位。对同一例子提交新的
  锯齿图。

  \item 对矩阵
  
  $$
  A=\operatorname{diag}(15:-1:1)+\operatorname{ones}(15,15)
  $$
  
  重新运行程序，并分别生成使用移位和不使用移位的两幅锯齿图。讨论在这个
  矩阵和前一个矩阵上观察到的收敛速率：它是线性的、超线性的、二次的、三次的，
  还是其他类型？谈论“每个特征值需要多少次 QR 迭代”有意义吗？
\end{enumerate}

\end{document}
